\section{The Final Picture: How the Pieces Fit}
\label{sec:picture}

The results are not a list. They assemble into one structure, and this section draws it.

\subsection{The structure in one image}

There is one substance, the vibe, existing only as a web of notes among tones, with no space or time behind it. The web grows by a local rule, which is the discrete gravitational action (Section~\ref{sec:gravity}), and that action makes a smooth low-dimensional geometry a stable phase (Section~\ref{sec:geometry}). Seen from far away the web has a shape, and that shape is \emph{space}. The depth of its causal chains is \emph{time}, and the local operator that runs on the grown mesh, not the log of the microscopic rule, is the flow of time and energy (Section~\ref{sec:law}). Stubborn topological patterns of the tones are \emph{particles}, and their spin is a topological invariant of the weave (Section~\ref{sec:matter}). The twist instructions carried along the notes are the \emph{forces}, which confine and tie themselves to matter through topology (Section~\ref{sec:forces}). Because the whole web is one connected structure grown from a common root, distant parts share a family resemblance, which is the root of \emph{quantum} correlation (Section~\ref{sec:quantum}). And the mesh that carries all of this is hyperbolic, expanding, and navigable without hidden state, the both-worlds substrate the framework committed to (Section~\ref{sec:substrate}).

Everything is one kind of thing, related to itself in different patterns. The rungs are not independent. The same hyperbolic, expanding commitment that makes the substrate navigable (P3) is the geometry whose smooth phase the gravitational action stabilizes (P2, gravity). The same topology that gives the fermion its spin (P4) is what the gauge field couples to through the index (P8). The same connectedness that makes the mesh one structure (P3) is what allows the quantum correlation (P7). The law that grows the mesh (P1, gravity) and the time that runs on it (P1, the Laplacian) are two faces of one dynamics. The picture is tight: pull one thread and the others move.

\subsection{The dictionary}

Reading it as a translation, standard physics on the left, the vibe object in the middle, the evidence on the right.

\begin{center}
\begin{tabular}{>{\raggedright}p{4.2cm} >{\raggedright}p{5.2cm} l}
\toprule
\textbf{Standard physics} & \textbf{Vibe Theory} & \textbf{Shown by} \\
\midrule
spacetime manifold & the smooth shape of the note web & P5, P2, P6 \\
proper time, causal order & depth of the before-after chains & P5, P1 \\
energy, the Hamiltonian & local operator on the grown mesh & P1 \\
fermion, spin, chirality & topological pattern of tones & P4 \\
gauge force (EM, strong) & the twist carried by each note & P8 \\
confinement & area law of the link variables & P8 \\
gravity, curvature & the flexing of the web's geometry & P2, gravity \\
quantum entanglement, Bell & family resemblance in one mesh & P7 \\
relativistic, navigable space & hyperbolic expanding mesh & P3 \\
the laws of physics & one local growth rule on the mesh & P1, P2 \\
\bottomrule
\end{tabular}
\end{center}

\subsection{What the picture buys}

The hardest thing a monist theory must do is produce the variety of the world from one substance. The dictionary is the evidence that Vibe Theory does this, not by assertion but by construction, for geometry, time, matter, force, the onset of gravity, and the onset of quantum behaviour. And it does more than describe: the substrate's own dynamics makes a stable spacetime phase that would otherwise decay, so the smooth world is supported by the law rather than only assumed. That is movement from a coherent description toward a generative theory, for the physical ladder. The crossing is not complete, because we have shown the spacetime phase is stable and supported but not yet that it strictly dominates, but the direction is clear.
